% Terminology and editorial conventions
% ========================================
%
% Gabriel's 1962 French coinages do not all survive into modern English, and
% several of them were superseded by Grothendieck's own usage within a few
% years.  The policy for this translation is:
%
%   1. Where standard English mathematical usage exists, use it.  A reader
%      should be able to find the result in the modern literature under the
%      word used here.
%   2. Where the original term is itself the point -- because it is a coinage
%      Gabriel is introducing, or because it names a notion that later
%      disappeared -- keep an accurate English rendering and give the French
%      original in parentheses at its first occurrence, and thereafter in the
%      terminology table below.
%   3. Sentences are translated, not paraphrased.  Definitions keep their
%      original form (data / axioms as bullet lists) so that the translation
%      can be checked against the original clause by clause.
%
% The table below is the single place where these decisions are recorded.

\begin{table}[htbp]
  \centering
  \small
  \begin{tabular}{p{0.34\textwidth} p{0.28\textwidth} p{0.28\textwidth}}
    \hline
    \textbf{French (original)} & \textbf{English (this translation)} & \textbf{Modern usage}\\[2pt]
    \hline
    univers                     & Grothendieck universe       & Grothendieck universe\\
    ensemble                    & set                          & set\\
    catégorie                   & category                     & category\\
    objet                       & object                       & object\\
    morphisme                   & morphism                     & morphism\\
    morphisme identique         & identity morphism            & identity map\\
    application de composition & composition                  & composition\\
    morphisme functoriel        & natural transformation       & natural transformation\\
    isomorphisme                & isomorphism                 & isomorphism\\
    fondeur / foncteur          & functor                      & functor\\
    Yoncteur                    & Yoneda embedding             & Yoneda embedding\\
    source / but                & source / target              & source / target\\
    mono\-morphisme             & monomorphism                 & monomorphism\\
    épi\-morphisme              & epimorphism                  & epimorphism\\
    sous-objet                  & subobject                    & subobject\\
    quotient                    & quotient                     & quotient\\
    sous-catégorie              & subcategory                 & subcategory\\
    sous-catégorie pleine       & full subcategory             & full subcategory\\
    catégorie duale             & dual category                & dual category\\
    catégorie produit           & product category             & product category\\
    indice                      & index set                    & indexing family\\
    système inductif            & inductive system             & directed system\\
    limite inductive            & inductive limit              & direct limit\\
    système projectif           & projective system            & inverse system\\
    limite projective            & projective limit             & inverse limit\\
    générateur                  & generator                    & generator\\
    générateurs                 & generators                   & generators\\
    ajouté                      & coproduct                    & coproduct\\
    produit                     & product                     & product\\
    \textsc{Ens}                & $\mathbf{Ens}$ (sets)        & $\mathbf{Set}$\\
    \textsc{Ab}                 & $\mathbf{Ab}$ (abelian groups) & $\mathbf{Abel}$\\
    noethérien (à droite)       & (right) noetherian           & noetherian\\
    noethérien                  & noetherian                   & noetherian\\
    localement noethérien       & locally noetherian          & locally noetherian\\
    dimension de Krull          & Krull dimension             & Krull dimension\\
    module pseudo-compact       & pseudocompact module         & pseudocompact module\\
    enveloppe injective         & injective envelope           & injective envelope\\
    foncteur exact à gauche    & left exact functor           & left exact functor\\
    catégorie localisante       & localizing subcategory      & localizing subcategory\\
    catégorie quotient          & quotient category            & quotient category\\
    langage modulo $\mathcal{C}$& language modulo $\mathcal{C}$ & Serre's $\bmod \mathcal{C}$\\
    recollement                 & gluing                       & gluing / recollement\\
    préschéma                   & prescheme                    & scheme / prescheme\\
    faisceau quasi cohérent     & quasi-coherent sheaf        & quasi-coherent sheaf\\
    \hline
  \end{tabular}
  \caption{Translation of the terminology.  Page references are to the
  original French edition, \emph{Bulletin de la Société Mathématique de
  France} \textbf{90} (1962), 323--448.}
  \label{tab:terminology}
\end{table}